% Nouveau EGA V (cote 25): archive cote 25, site page 162 (PDF page 163).
% English translation of the checked French annotated master.
% AI-translated by GPT-6 Astra (OpenAI, ChatGPT Pro mode).
\scanpage{25/162}
\mspage{\hand{\unsure{5}}}

$\hand{\geqslant}$ 2, we shall say that x \hand{$\in$} X is an elementary singular point of multiplicity
n if its local ring $\underline{O}_{X,x}$ is elementary of multiplicity n. (Point out
in a remark that this \var{\xout{notion}}{}{typewriter strike-out}{ega25_25_162-v0001} terminology agrees with the general notion
of multiplicity due to Samuel). For n=2, we shall speak in particular
of an elementary quadrat\unsure{i|t}c singularity (or ordinary, in the classical
terminology). Similarly, in accordance with general usage, one int\unsure{r|o}oduces
the ``geometric'' variants for X locally of finite type over a field:
x \var{\inserted[typed]{$\hand{\in} X$}}{$\hand{\in} X$}{insertion}{ega25_25_162-v0002} is called a geometrically elementary singularity of multiplicity n if,
for every (or, equivalently\var{\over{\unsure{;}}{\unsure{,}}}{\unsure{,}}{writing over}{ega25_25_162-v0003} one) extension K of k, and \unsure{every}
(or one) point$^{z}$ of $X_{K}$ above x, rational over K, z is an
elementary singularity of multiplicity n.

15.3. ``Generalize'' SUP prop.8 to the case of multiplicity n.

15.4. ``Generalize'' SUP prop.9 and cor.10. to the case of multiplicity n.
(These are not really generalizations, but variants..)

15.5. Introduce the notion of an elementary singular zero of multiplicity
n (for n=2, an elementary \unsure{quadra}\unsure{tic} si\unsure{n|g}g\unsure{u|l}\unsure{l|u}ar zero, or an ordinary one) of a
section $\hand{\phi}$ of $\underline{O}_{X}$, or more generally of a locally free Module, on
a locally noetherian X, and the corresponding geometric notion
(over a base field k).

15.6. Let X be a subprescheme of a regular loc noeth prescheme X', and
let x \hand{$\in$}X. For \unsure{x|X} to be an elementary singularity of multiplicity
n of X, it is necessary and sufficient that there be an open neighborhood U of x, and an
$\underline{O}_{U}$-regular sequence ($\hand{\phi}_{1},..,\hand{\phi}_{d}$) such that $X|U = V(\hand{\phi}_{1},..,\hand{\phi}_{d})$ and
such that x is an elementary singular zero of multiplicity n of ($\hand{\phi}_{1}$,..).
In particular, if A is a local noetherian ring that is elementary
\var{\margin{$\Vert$}}{\placedmargin{$\Vert$}}{marginal note placed}{ega25_25_162-v0006} of multiplicity n, then A is a ``complete intersection'' ring; moreover,
one sees by 15\unsure{.|1}4. that one can find, in a neighborhood of x,

\mspage{\hand{162}}
\tnote{25-162-1}{The raised z after point remains in its printed position. Every retains the French initial t/d doubt in tout. In the parenthetical term of15.5, the two linked parts of quadratic retain the separate i/t and u/q doubts in quadratique; the three doubts in singular retain the source n/g, u/l and l/u. The overwritten doubtful g/q there becomes g in French without changing English singular, so that variant is omitted and the old \unsure{g|q} is retained here.}
\tnote{25-162-2}{The canceled feminine e before voisinage has no English effect on an open neighborhood. The canceled fragment \unsure{ieu|leu} within d’ailleurs has no English counterpart. Loc, noeth, prop. and cor. mean locally, noetherian, proposition and corollary; il f et s is rendered it is necessary and sufficient. The reference15.4 retains its doubtful separator.}
